\section{Key concepts}
\label{sec:concepts}

\subsection{Image and caption retrieval}
An \emph{embedding} is a list of numbers that represents an image or a text caption. Embeddings from a shared model can be compared to find items with similar representations. In this project, the image encoder is used to retrieve training images similar to a held-out query image. The caption index stores vectors for training captions, but a caption-generation pipeline that uses these results has not been implemented.

\subsection{What similarity does and does not mean}
The project uses CLIP, a vision-language model trained to align image and text representations \citep{radford2021clip,clipmodelcard}. A high similarity score means that two vectors are close under the chosen representation. It does not prove that the retrieved image is an exact match, that its caption applies to the query, or that a future generated caption is factually correct.

FAISS is a library for efficient similarity search over vectors \citep{douze2024faiss,faissdocs}. The project uses its exact flat inner-product index. Because the stored vectors are L2-normalized, inner product gives cosine similarity. This index provides a straightforward reference search, not an approximate search optimization.

\subsection{Phase boundaries}
Phase 1 ends with saved splits, analysis, embeddings, indexes, and retrieval checks. Phase 2 must add a generator, define retrieval conditions and prompts, run controlled generation, and evaluate the outputs. The two phases should not be confused when interpreting the results in this report.
